
quelconques \(c_{h,k,\lambda}\) et \(c_{h,k,\mu}\) sont conjuguées, autrement dit il existe une matrice inversible \(a=(a_{ij})\) \((i,j=1,\ldots,n)\) telle que l'on ait
\[
 ac_{h,k,\lambda}=c_{h,k,\mu}a.
\]
En effet, la matrice \((a_{ij})\) avec \(a_{hh}=1/\mu\), \(a_{kk}=1/\lambda\), et du reste \(a_{ij}=\delta_{ij}\), répond à la question (rappelons que \(\lambda\) et \(\mu\) sont par hypothèse inversibles et \(h\ne k\)). Donc pour un homomorphisme quelconque \(\beta\) de \(\operatorname{G\ell}(n)(U)\) dans un faisceau de groupes abéliens \(H\), on a
\[
 \beta(c_{h,k,\lambda})=\beta(c_{h,k,\mu}).
\]
Si \(\lambda+\mu\) est inversible, on tire de ceci et de (6.3.1) que \(\beta(c_{h,k,\lambda})=0\) pour tout \(h,\beta\), de sorte que \(c_{h,k,\lambda}\) est un produit de commutateurs. Si le corps résiduel \(k\) de \(y\) n'est pas le corps premier \(\underline{F}_2=\{0,1\}\), alors on peut trouver pour tout \(\lambda\in A\) inversible un voisinage zariskien \(U=\Spec(B)\) de \(y\) et un élément inversible \(\mu\in B\) tel que \(\lambda+\mu\in B\) soit inversible. Il en découle que \(c_{h,k,\lambda}\) est un produit de commutateurs sur \(U\).

Si \(k=\underline{F}_2\), alors \(V=\Spec(A[T]/(T^2+T+1))\) est un voisinage étale de \(y\), dont le corps résiduel est différent de \(\underline{F}_2\). On a gagné.

\subsection*{6.5. Remarques}
\textbf{6.5.1.} Lorsque les corps résiduels de \(S\) sont tous non isomorphes à \(\underline{F}_2\), la démonstration précédente montre que 6.2.2. reste valable pour le topos annelé des faisceaux \emph{zariskiens} sur \(S\). En revanche, déjà pour \(S=\Spec(\underline{F}_2)\), \(\operatorname{S\ell}(2,\underline{F}_2)=\operatorname{G\ell}(2,\underline{F}_2)\) est différent du groupe des commutateurs de \(\operatorname{G\ell}(2,\underline{F}_2)\). L'ordre de ce dernier est \(3\) tandis que l'ordre de \(\operatorname{S\ell}(2,\underline{F}_2)\) est \(6\) (cf. [1]).

\textbf{6.5.2.} Sur un topos localement annelé quelconque \(X\) il n'est pas vrai en général que toute représentation \(G_m\to G_m\) soit définie par une
